%! AP Statistics — Unit 9, Lesson 9.6 — Homework (Answer Key)
\documentclass[10pt]{article}
\usepackage{apstats-article}
\usepackage{apstats-key}
\usepackage{grading}
\gradingassignment{apstats-u09-l06-hw}
\gradingcourse{advanced_statistics}

\begin{document}

\pageheader{Unit 9, Lesson 9.6}{Homework --- Answer Key}

\vspace{0.1in}

\textbf{Directions:} For each scenario, answer all three parts.
\begin{itemize}[noitemsep]
  \item[(1)] Name the inference procedure.
  \item[(2)] State the parameter being estimated (for a CI) or the hypotheses $H_0$ and $H_a$ (for a test).
  \item[(3)] List the key conditions that must be checked for this procedure.
\end{itemize}

\boxguard[15]
\begin{notesbox}{Problem 1}
\par\noindent\begin{minipage}{\linewidth}
A random sample of 45 high school seniors records their SAT scores. The researcher wants
to estimate the mean SAT score for all seniors at this school with a 95\% confidence interval.

\vspace{0.08in}
(1) Procedure: \ans{1-sample $t$-interval for a mean}\gradedkeyline{\gradedblank[id=hw1a, type=rubric, points=1, display={\bfseries\boldmath 1-sample $t$-interval for $\mu$}]{\linewidth+6pt}{\gradedlinesheight{9}}}
\end{minipage}

\vspace{0.08in}
\par\noindent\begin{minipage}{\linewidth}
(2) Parameter being estimated: \ans{$\mu$ = the true mean SAT score of all seniors at this school}\gradedkeyline{\gradedblank[id=hw1b, type=rubric, points=1, display={\bfseries\boldmath $\mu$ = true mean SAT score, all seniors here}]{\linewidth+6pt}{\gradedlinesheight{9}}}
\end{minipage}

\vspace{0.08in}
\par\noindent\begin{minipage}{\linewidth}
(3) Conditions: \ans{Random sample; Normal/Large Sample ($n = 45 \ge 30$, so the condition is met); 10\% rule: $45 \le 10\%$ of all seniors at the school.}\gradedkeyline{\gradedblank[id=hw1c, type=rubric, points=2, display={\bfseries\boldmath Random; $n = 45 \ge 30$; $45 \le 10\%$ of seniors.}]{\linewidth+6pt}{\gradedlinesheight{1}}}
\end{minipage}
\end{notesbox}

\newpage
\boxguard[15]
\begin{notesbox}{Problem 2}
\par\noindent\begin{minipage}{\linewidth}
A school randomly surveys 120 sophomores about how they get to school (walk, bus, car, or
bike). The researcher tests whether the four options are equally popular.

\vspace{0.08in}
(1) Procedure: \ans{Chi-square goodness of fit test}\gradedkeyline{\gradedblank[id=hw2a, type=rubric, points=1, display={\bfseries\boldmath $\chi^2$ goodness-of-fit test}]{\linewidth+6pt}{\gradedlinesheight{9}}}
\end{minipage}

\vspace{0.08in}
\par\noindent\begin{minipage}{\linewidth}
(2) $H_0$: \ans{$p_{\text{walk}} = p_{\text{bus}} = p_{\text{car}} = p_{\text{bike}} = 0.25$}\gradedkeyline{\gradedblank[id=hw2b1, type=rubric, points=1, display={\bfseries\boldmath All four $p = 0.25$}]{\linewidth+6pt}{\gradedlinesheight{9}}}
\end{minipage}

\vspace{0.08in}
\par\noindent\begin{minipage}{\linewidth}
$H_a$: \ans{At least one of the proportions differs from 0.25.}\gradedkeyline{\gradedblank[id=hw2b2, type=rubric, points=1, display={\bfseries\boldmath Some $p \ne 0.25$}]{\linewidth+6pt}{\gradedlinesheight{9}}}
\end{minipage}

\tcbbreak

\vspace{0.08in}
\par\noindent\begin{minipage}{\linewidth}
(3) Conditions: \ans{Random sample; Large Counts: all expected counts $= 120 \times 0.25 = 30 \ge 5$.}\gradedkeyline{\gradedblank[id=hw2c, type=rubric, points=2, display={\bfseries\boldmath Random; all expected $= 30 \ge 5$.}]{\linewidth+6pt}{\gradedlinesheight{1}}}
\end{minipage}
\end{notesbox}

\boxguard[15]
\begin{notesbox}{Problem 3}
\par\noindent\begin{minipage}{\linewidth}
A researcher randomly selects 80 adults in a city and classifies each as a smoker or
non-smoker. She tests whether the proportion of smokers in the city differs from the
national rate of 14\%.

\vspace{0.08in}
(1) Procedure: \ans{1-proportion $z$-test}\gradedkeyline{\gradedblank[id=hw3a, type=rubric, points=1, display={\bfseries\boldmath 1-proportion $z$-test}]{\linewidth+6pt}{\gradedlinesheight{9}}}
\end{minipage}

\vspace{0.08in}
\par\noindent\begin{minipage}{\linewidth}
(2) $H_0$: \ans{$p = 0.14$}\gradedkeyline{\gradedblank[id=hw3b1, type=rubric, points=1, display={\bfseries\boldmath $p = 0.14$}]{\linewidth+6pt}{\gradedlinesheight{9}}}
\end{minipage}

\tcbbreak

\vspace{0.08in}
\par\noindent\begin{minipage}{\linewidth}
$H_a$: \ans{$p \ne 0.14$}\gradedkeyline{\gradedblank[id=hw3b2, type=rubric, points=1, display={\bfseries\boldmath $p \ne 0.14$}]{\linewidth+6pt}{\gradedlinesheight{9}}}
\end{minipage}

\vspace{0.08in}
\par\noindent\begin{minipage}{\linewidth}
(3) Conditions: \ans{Random sample; 10\% rule: $80 \le 10\%$ of all adults in the city; Large Counts: $np_0 = 80(0.14) = 11.2 \ge 10$ and $n(1-p_0) = 80(0.86) = 68.8 \ge 10$.}\gradedkeyline{\gradedblank[id=hw3c, type=rubric, points=2, display={\bfseries\boldmath Random; $80 \le 10\%$; $11.2, 68.8 \ge 10$.}]{\linewidth+6pt}{\gradedlinesheight{1}}}
\end{minipage}
\end{notesbox}

\boxguard[15]
\begin{notesbox}{Problem 4}
\par\noindent\begin{minipage}{\linewidth}
A pediatrician follows 60 children for 5 years, recording average daily screen time
(hours/day, $x$) and reading proficiency score ($y$). She wants to test whether screen
time is linearly related to reading proficiency.

\vspace{0.08in}
(1) Procedure: \ans{$t$-test for slope}\gradedkeyline{\gradedblank[id=hw4a, type=rubric, points=1, display={\bfseries\boldmath $t$-test for slope}]{\linewidth+6pt}{\gradedlinesheight{9}}}
\end{minipage}

\tcbbreak

\vspace{0.08in}
\par\noindent\begin{minipage}{\linewidth}
(2) $H_0$: \ans{$\beta = 0$ (screen time is not a useful linear predictor of reading proficiency)}\gradedkeyline{\gradedblank[id=hw4b1, type=rubric, points=1, display={\bfseries\boldmath $\beta = 0$}]{\linewidth+6pt}{\gradedlinesheight{9}}}
\end{minipage}

\vspace{0.08in}
\par\noindent\begin{minipage}{\linewidth}
$H_a$: \ans{$\beta \ne 0$}\gradedkeyline{\gradedblank[id=hw4b2, type=rubric, points=1, display={\bfseries\boldmath $\beta \ne 0$}]{\linewidth+6pt}{\gradedlinesheight{9}}}
\end{minipage}

\vspace{0.08in}
\par\noindent\begin{minipage}{\linewidth}
(3) Conditions: \ans{LINER --- Linear (residual plot shows no pattern); Independent observations; Normal residuals (check histogram/boxplot of residuals); Equal variance (residual plot shows constant spread); Random sample.}\gradedkeyline{\gradedblank[id=hw4c, type=rubric, points=2, display={\bfseries\boldmath LINER: Linear, Independent, Normal residuals, Equal SD, Random.}]{\linewidth+6pt}{\gradedlinesheight{1}}}
\end{minipage}
\end{notesbox}

\newpage
\boxguard[15]
\begin{notesbox}{Problem 5}
\par\noindent\begin{minipage}{\linewidth}
A study records whether a randomly selected sample of 200 adults (100 from urban areas,
100 from rural areas) experienced anxiety in the past year (yes/no). The researcher
tests whether anxiety rates differ between urban and rural adults.

\vspace{0.08in}
(1) Procedure: \ans{2-proportion $z$-test}\gradedkeyline{\gradedblank[id=hw5a, type=rubric, points=1, display={\bfseries\boldmath 2-proportion $z$-test}]{\linewidth+6pt}{\gradedlinesheight{9}}}
\end{minipage}

\vspace{0.08in}
\par\noindent\begin{minipage}{\linewidth}
(2) $H_0$: \ans{$p_{\text{urban}} = p_{\text{rural}}$ (anxiety rates are the same in both groups)}\gradedkeyline{\gradedblank[id=hw5b1, type=rubric, points=1, display={\bfseries\boldmath $p_{\text{urban}} = p_{\text{rural}}$}]{\linewidth+6pt}{\gradedlinesheight{9}}}
\end{minipage}

\vspace{0.08in}
\par\noindent\begin{minipage}{\linewidth}
$H_a$: \ans{$p_{\text{urban}} \ne p_{\text{rural}}$}\gradedkeyline{\gradedblank[id=hw5b2, type=rubric, points=1, display={\bfseries\boldmath $p_{\text{urban}} \ne p_{\text{rural}}$}]{\linewidth+6pt}{\gradedlinesheight{9}}}
\end{minipage}

\tcbbreak

\vspace{0.08in}
\par\noindent\begin{minipage}{\linewidth}
(3) Conditions: \ans{Random sample from each group; 10\% rule for each group; Large Counts: $n\hat{p} \ge 10$ and $n(1-\hat{p}) \ge 10$ for each group (using the combined $\hat{p}$).}\gradedkeyline{\gradedblank[id=hw5c, type=rubric, points=2, display={\bfseries\boldmath Random each group; 10\% each; Large Counts (pooled $\hat p$) $\ge 10$.}]{\linewidth+6pt}{\gradedlinesheight{1}}}
\end{minipage}
\end{notesbox}

\end{document}
